\documentclass[11pt]{article}
\usepackage[margin=1in]{geometry}
\usepackage{amsmath, amssymb, amsthm, mathtools}
\usepackage{enumitem}
\usepackage{booktabs}
\usepackage{array}

\theoremstyle{plain}
\newtheorem{theorem}{Theorem}[section]
\newtheorem{proposition}[theorem]{Proposition}
\newtheorem{lemma}[theorem]{Lemma}
\newtheorem{corollary}[theorem]{Corollary}
\theoremstyle{definition}
\newtheorem{definition}[theorem]{Definition}
\newtheorem{remark}[theorem]{Remark}

\newcommand{\cA}{\mathcal{A}}
\newcommand{\cB}{\mathcal{B}}
\newcommand{\cZ}{\mathcal{Z}}
\newcommand{\cV}{\mathcal{V}}
\newcommand{\cN}{\mathcal{N}}
\DeclareMathOperator{\rank}{rank}

\title{Rodina's Hidden-Zero Theorem: The $n=6$ Base Case\\[0.3em]
\large Direct proof that cyclic $1$-zeros uniquely determine $\cA_6^{\mathrm{tree}}$}
\author{Jonathan Valenzuela}
\date{\today}

\begin{document}
\maketitle

\begin{abstract}
We prove by direct linear-algebraic computation that the $84$-dimensional simple-pole
ansatz at $n=6$ is reduced to a $1$-dimensional nullspace (spanned by the $\phi^3$
tree amplitude) upon imposing all six cyclic $1$-zero conditions. The proof is
structured leg-by-leg, with the exact rank increase at each step exhibited. Combined
with the $n=5$ base case (Part~I), this provides two fully rigorous base cases for
the general hidden-zero uniqueness conjecture. The computation is verified
symbolically and the key structural features---including the precise number of
independent $1$-zeros and the intermediate non-local survivors---are identified.
\end{abstract}

%======================================================================
\section{Setup at $n=6$}
%======================================================================

\subsection{Kinematic space}

At $n=6$ the planar Mandelstam variables are
\[
X_{13},\; X_{14},\; X_{15},\; X_{24},\; X_{25},\; X_{26},\; X_{35},\; X_{36},\; X_{46},
\]
forming a $9$-dimensional kinematic space $\cV_6$.

\subsection{Ansatz space}

\begin{definition}
The simple-pole ansatz space $\cB_6$ consists of all rational functions
\[
B \;=\; \sum_{\alpha} c_\alpha \cdot \frac{1}{X_{a_1} X_{a_2} X_{a_3}},
\]
where $\alpha = (a_1, a_2, a_3)$ ranges over all $\binom{9}{3} = 84$ unordered
triples of distinct planar Mandelstam variables. Each $B \in \cB_6$ is homogeneous
of mass dimension $-3$.
\end{definition}

\subsection{The six cyclic $1$-zero loci}

Each locus $\cZ_k$ is defined by $n-3 = 3$ linear constraints. The explicit
substitution rules are given in Table~\ref{tab:loci6}.

\begin{table}[ht]
\centering
\small
\begin{tabular}{c|l|l}
\toprule
$k$ & Substitution (eliminated $\to$ free) & Free variables \\
\midrule
$1$ & $X_{13} = -X_{26},\; X_{14} = X_{24}-X_{26},\; X_{15} = X_{25}-X_{26}$
    & $X_{24}, X_{25}, X_{26}, X_{35}, X_{36}, X_{46}$ \\
$2$ & $X_{13} = X_{36}-X_{26},\; X_{24} = X_{26}-X_{36},\; X_{25} = X_{26}+X_{35}-X_{36}$
    & $X_{14}, X_{15}, X_{26}, X_{35}, X_{36}, X_{46}$ \\
$3$ & $X_{13} = X_{14}+X_{36}-X_{46},\; X_{24} = X_{46}-X_{36},\; X_{35} = X_{36}-X_{46}$
    & $X_{14}, X_{15}, X_{25}, X_{26}, X_{36}, X_{46}$ \\
$4$ & $X_{14} = X_{15}+X_{46},\; X_{24} = X_{25}+X_{46},\; X_{35} = -X_{46}$
    & $X_{13}, X_{15}, X_{25}, X_{26}, X_{36}, X_{46}$ \\
$5$ & $X_{15} = -X_{46},\; X_{25} = X_{26}-X_{46},\; X_{35} = X_{36}-X_{46}$
    & $X_{13}, X_{14}, X_{24}, X_{26}, X_{36}, X_{46}$ \\
$6$ & $X_{13} = X_{36}-X_{26},\; X_{14} = X_{46}-X_{26},\; X_{15} = -X_{26}$
    & $X_{24}, X_{25}, X_{26}, X_{35}, X_{36}, X_{46}$ \\
\bottomrule
\end{tabular}
\caption{The six cyclic $1$-zero loci at $n=6$.}
\label{tab:loci6}
\end{table}

\subsection{The tree amplitude}

The ordered $\phi^3$ tree amplitude at $n=6$ is the sum over all $C_4 = 14$
non-crossing triangulations of the convex hexagon:
\begin{equation}
\cA_6^{\mathrm{tree}} \;=\; \sum_{\text{14 triangulations}} \frac{1}{X_{a_1} X_{a_2} X_{a_3}}.
\label{eq:tree6}
\end{equation}
The $14$ triangulations (as triples of diagonals) are listed in Table~\ref{tab:tri6}.

\begin{table}[ht]
\centering
\small
\begin{tabular}{llll}
\toprule
$(X_{13}, X_{14}, X_{15})$ & $(X_{13}, X_{14}, X_{46})$ &
$(X_{13}, X_{15}, X_{35})$ & $(X_{13}, X_{35}, X_{36})$ \\
$(X_{13}, X_{36}, X_{46})$ & $(X_{14}, X_{15}, X_{24})$ &
$(X_{14}, X_{24}, X_{46})$ & $(X_{15}, X_{24}, X_{25})$ \\
$(X_{15}, X_{25}, X_{35})$ & $(X_{24}, X_{25}, X_{26})$ &
$(X_{24}, X_{26}, X_{46})$ & $(X_{25}, X_{26}, X_{35})$ \\
$(X_{26}, X_{35}, X_{36})$ & $(X_{26}, X_{36}, X_{46})$ & & \\
\bottomrule
\end{tabular}
\caption{The $14$ non-crossing triangulations of the hexagon.}
\label{tab:tri6}
\end{table}

%======================================================================
\section{The $1$-zero conditions as a linear system}
%======================================================================

The procedure is identical to Part~I at $n=5$: for each leg $k$, substitute the
$\cZ_k$ relations into the generic ansatz $B = \sum_{i=1}^{84} c_i\, t_i$ and clear
denominators. The numerator, as a polynomial in the $6$ free variables on $\cZ_k$,
must vanish identically. Setting each monomial coefficient to zero yields a system
of linear equations in the $84$ coefficients $c_1, \ldots, c_{84}$.

We process legs sequentially, tracking the cumulative rank and nullspace dimension.

\subsection{Leg-by-leg rank analysis}

\begin{theorem}\label{thm:n6-rank}
The rank of the cumulative constraint matrix evolves as follows:

\begin{center}
\begin{tabular}{lccc}
\toprule
Legs imposed & Raw equations & Cumulative rank (of $84$) & Nullspace dim \\
\midrule
Leg $1$ only      & $63$  & $51$ & $33$ \\
Legs $1, 2$       & $154$ & $76$ & $8$  \\
Legs $1, 2, 3$    & $245$ & $81$ & $3$  \\
Legs $1, 2, 3, 4$ & $308$ & $82$ & $2$  \\
Legs $1, \ldots, 5$ & $371$ & $83$ & $1$ \\
All $6$ legs      & $434$ & $83$ & $1$  \\
\bottomrule
\end{tabular}
\end{center}

The final nullspace is $1$-dimensional and spanned by $\cA_6^{\mathrm{tree}}$.
\end{theorem}

\begin{proof}
Verified by exact symbolic computation in \texttt{sympy}. The computation
substitutes each $\cZ_k$ into the $84$-term ansatz, clears denominators to obtain
a polynomial in the $6$ free variables, extracts all monomial coefficient equations,
row-reduces the accumulated system, and reports the rank. The nullspace vector
has exactly $14$ nonzero entries, all equal, corresponding to the $14$ triangulations
of the hexagon.
\end{proof}

\subsection{Structural observations}

\begin{corollary}[Independence count at $n=6$]
Exactly $5$ of the $6$ cyclic $1$-zeros are independent as constraints on $\cB_6$.
Legs $1$ through $5$ suffice; leg $6$ is redundant.
\end{corollary}

\begin{remark}[Comparison with $n=5$]
At $n=5$, $3$ of $5$ legs suffice (legs $4, 5$ redundant).
At $n=6$, $5$ of $6$ legs suffice (leg $6$ redundant).
The number of independent legs is $n-1$ at both $n=5$ and $n=6$.
\end{remark}

\begin{remark}[Contribution of each leg]
The contributions are highly non-uniform: leg~$1$ contributes $51$ independent
equations (killing $60\%$ of the ansatz), leg~$2$ adds $25$ more ($30\%$), and
legs $3$--$5$ contribute $5 + 1 + 1 = 7$ together ($8\%$). The first two legs
do the overwhelming majority of the work, reducing the nullspace from $84$ to $8$.
The remaining legs trim the $8$-dimensional survivor space down to $1$.
\end{remark}

%======================================================================
\section{The mechanism: grading by free variables and limits}
%======================================================================

We now describe the structural mechanism behind leg~$1$'s $51$ equations,
following the argument of Part~I and the author's handwritten notes.

\subsection{Free-variable grading on $\cZ_1$}

On $\cZ_1$, the $6$ free variables split into two groups:
\begin{itemize}[nosep]
\item \textbf{Fully free:} $X_{35}, X_{36}, X_{46}$ --- these do not appear in any
      substitution expression on the right-hand side of Table~\ref{tab:loci6}.
\item \textbf{Constrained-free:} $X_{24}, X_{25}, X_{26}$ --- these appear in the
      substitution expressions ($X_{13} = -X_{26}$, $X_{14} = X_{24} - X_{26}$,
      $X_{15} = X_{25} - X_{26}$).
\end{itemize}

Since $X_{35}, X_{36}, X_{46}$ are genuinely free parameters of $\cZ_1$ and do not
enter the constraint equations, any rational function that must vanish identically
on $\cZ_1$ can be \emph{graded} by the powers of these variables appearing in
denominators, and each grade must vanish independently (by linear independence of
Laurent monomials in the free variables).

\subsection{What the grading kills}

For each of the $84$ basis terms $1/(X_a X_b X_c)$, count how many of $\{a, b, c\}$
lie in the fully-free set $\{X_{35}, X_{36}, X_{46}\}$. Call this the \emph{free grade}
$g \in \{0, 1, 2, 3\}$. The distribution is:

\begin{center}
\begin{tabular}{ccc}
\toprule
Free grade $g$ & Number of terms & Interpretation \\
\midrule
$3$ & $1$  & $1/(X_{35} X_{36} X_{46})$ \\
$2$ & $18$ & Two free + one other \\
$1$ & $45$ & One free + two others \\
$0$ & $20$ & No free variables \\
\bottomrule
\end{tabular}
\end{center}

The grade-$3$ term and grade-$2$ terms are the ``easiest'' to constrain: grade~$3$
is a single coefficient that must vanish; grade~$2$ terms, upon restriction to
$\cZ_1$, become sums $\sum_j c_j / x_j$ with the $x_j$ depending on the
constrained-free and eliminated variables after substitution, and the remaining
fully-free variables can be used to isolate individual coefficients via limits.

The grade-$1$ and grade-$0$ terms are more intricate: their vanishing involves
\emph{collective cancellations} between terms whose denominators are related by the
constraint substitutions (e.g., $X_{13} = -X_{26}$ pairs terms containing $X_{13}$
with terms containing $X_{26}$). These cancellations produce the coefficient
\emph{equalities} (rather than vanishings) that identify the tree amplitude.

\subsection{Analogy with $n=5$}

At $n=5$, $\cZ_1$ has one fully-free variable ($X_{35}$) and two constrained-free
variables ($X_{24}, X_{25}$). The grading by $X_{35}$ splits the $15$-term ansatz
into grades $0$ (10 terms), $1$ (4~terms), and $2$ (1~term). The limit argument
at $n=5$ killed grade~$2$ and constrained grade~$1$, leaving a $6$-dimensional
survivor; cycling through legs $2$ and $3$ then reduced to $1$ dimension.

At $n=6$ the same mechanism operates with three fully-free variables instead of one,
producing a finer grading and more powerful constraints from a single leg. This
accounts for leg~$1$ alone killing $51$ of $84$ dimensions at $n=6$ (vs.\ $9$ of
$15$ at $n=5$): both represent $\sim 60\%$ of the ansatz, achieved by the
free-variable grading.

%======================================================================
\section{The intermediate survivors}
%======================================================================

After leg~$1$ (nullspace dim $= 33$) and leg~$2$ (nullspace dim $= 8$), there is an
$8$-dimensional space of rational functions satisfying the first two $1$-zeros. This
space contains the $1$-dimensional tree-amplitude direction plus a $7$-dimensional
space of \emph{non-local survivors}---rational functions with poles in non-planar
channels that are not yet eliminated.

At $n=5$, the analogous intermediate survivor after two legs was $1$-dimensional
(the non-local function $\cN$ from Part~I, eq.~(8)). At $n=6$, the richer
structure (more variables, more crossing-chord triples) produces a larger
intermediate space, but the subsequent legs $3$--$5$ eliminate it completely.

The final nullspace after all legs is $1$-dimensional:
\begin{equation}
\ker\bigl(\text{all $6$ cyclic $1$-zeros on } \cB_6\bigr)
\;=\; \mathrm{span}\bigl\{\cA_6^{\mathrm{tree}}\bigr\}.
\end{equation}

%======================================================================
\section{The main theorem at $n=6$}
%======================================================================

\begin{theorem}\label{thm:n6}
Let $B \in \cB_6$ be a rational function of mass dimension $-3$ with at most simple
poles in the $9$ planar Mandelstam variables. If $B$ vanishes on each of the six
cyclic $1$-zero loci $\cZ_1, \ldots, \cZ_6$, then
\[
B \;=\; c \cdot \cA_6^{\mathrm{tree}}
\]
for some scalar $c$.
\end{theorem}

\begin{proof}
The $84$-dimensional ansatz $\cB_6$ is constrained by the six cyclic $1$-zero
conditions, which impose a total of $83$ independent linear equations on the
coefficients (Theorem~\ref{thm:n6-rank}). The solution space has dimension
$84 - 83 = 1$, and the unique generator is verified to be $\cA_6^{\mathrm{tree}}$
(all $14$ non-crossing triangulation coefficients equal, all others zero).
\end{proof}

\begin{corollary}
At $n=6$, locality (poles only in planar channels) and unitarity (correct
factorization at each pole) are consequences of the $1$-zero conditions, not
independent assumptions. Every non-planar pole (crossing-chord triple) is
eliminated by the $1$-zeros.
\end{corollary}

%======================================================================
\section{Combined status: $n = 5, 6, 7$}
%======================================================================

\begin{center}
\begin{tabular}{cccccc}
\toprule
$n$ & $\dim \cB_n$ & Rank & Nullspace & Method & Locality derived? \\
\midrule
$5$ & $15$   & $14$  & $1$ & Symbolic (exact) & Yes \\
$6$ & $84$   & $83$  & $1$ & Symbolic (exact) & Yes \\
$7$ & $1001$ & $1000$& $1$ & Numerical (SVD)  & Yes \\
\bottomrule
\end{tabular}
\end{center}

At each $n$, the cyclic $1$-zeros impose exactly $\dim \cB_n - 1$ independent
constraints on the full non-local ansatz, leaving a $1$-dimensional nullspace
spanned by the tree amplitude. Every crossing-chord and non-planar pole is
eliminated. The $n = 5$ and $n = 6$ results are proven by exact symbolic
computation; the $n = 7$ result is verified numerically with a singular-value
gap exceeding $13$ orders of magnitude.

\begin{remark}[Conjecture for general $n$]
We conjecture that for all $n \geq 5$,
\[
\dim \ker\bigl(\text{all $n$ cyclic $1$-zeros on } \cB_n\bigr) \;=\; 1,
\]
with the unique generator being $\cA_n^{\mathrm{tree}}$. The three verified base
cases ($n = 5, 6, 7$) and the structural uniformity of the free-variable grading
mechanism provide strong evidence for this conjecture.
\end{remark}

%======================================================================
\section*{Computational verification}
%======================================================================

All results are verified by the \texttt{sympy} scripts in
\texttt{computations/full\_ansatz\_n5.py} and
\texttt{computations/full\_ansatz\_n6.py}. The leg-by-leg rank computation at $n=6$
(Table in Theorem~\ref{thm:n6-rank}) was performed by sequentially substituting each
$\cZ_k$ into the $84$-term ansatz, clearing denominators, extracting monomial
coefficient equations, and row-reducing the accumulated system. The total computation
time is approximately $3$ minutes in \texttt{sympy} on a standard machine.

\end{document}
